\section{Read-level accuracy versus sample-level utility}
\label{sec:sample}

Read-level accuracy measures raw discriminative power, but practitioners
usually care about the community profile: the relative abundance of each taxon
in a sample. These two quantities can diverge, because independent per-read
errors partially cancel when aggregated. We therefore evaluate every model at
the sample level---relative-abundance estimation and binary detection---and
compare against Kraken\,2 on a database-coverage-matched pool.

\subsection{Per-read errors largely cancel in abundance estimates}

On synthetic species-balanced samples (9 samples $\times$ 10K reads from the
natural-distribution pool), a 67\%-accurate read classifier yields a
genus relative-abundance Pearson $r=0.993$ and Bray--Curtis dissimilarity
$0.098$ (Table~\ref{tab:sample}, Figure~\ref{fig:abund}). The overlapping-13-mer
model, with far higher read accuracy, reaches $r\!\approx\!1.000$. The ranking
by sample-level fidelity follows the ranking by read accuracy, but even the
mid-accuracy model is highly usable for abundance because misclassifications
are not systematically biased toward any one genus and thus average out. We
report $r$ as a range rather than a single value, because it is pool-dependent
(e.g.\ NT-v2 gives $r=0.928$ on the harder leftover pool); the genus-balanced
model, consistent with Section~\ref{sec:scaling}, is the clear outlier at
$r=0.862$.

We deliberately distinguish three levels of sample-level evaluation, of
increasing difficulty: (i) fixed-composition random partitions of a common read
pool (used here); (ii) sparse synthetic communities with variable composition;
and (iii) real or mock communities with authentic error profiles and incomplete
database coverage. The present results establish (i)---a stability test under a
fixed expected composition---and should not be read as evidence for (iii): they
show that per-read errors cancel when the community is held fixed, not that a
6-mer model can profile an arbitrary real sample. We therefore probe (iii) later
in this section on a real mock community (Section~\ref{sec:mock}); that probe is
preliminary, and a fuller multi-community evaluation remains required before
deployment (Section~\ref{sec:limitations}).

\begin{table}[t]
\centering
\caption{\textbf{Sample-level genus abundance} on the natural-distribution
pool (9$\times$10K reads). Higher read accuracy tracks better abundance
fidelity, but even a 67\%-accurate model attains $r=0.99$. Genus-balancing
(bottom row) is counterproductive.}
\label{tab:sample}
\begin{tabular*}{\columnwidth}{@{\extracolsep{\fill}}lccc@{}}
\toprule
Model & Read acc.\ & Pearson $r$ & Bray--Curtis \\
\midrule
MT 13-mer (250M)         & 98.7\% & 1.000 & 0.005 \\
MT 13-mer (50M)          & 87.4\% & 0.999 & 0.032 \\
NT-v2 6-mer (50M)        & 67.1\% & 0.993 & 0.098 \\
NT-v2 6-mer (250M warm)  & 67.3\% & 0.993 & 0.098 \\
MT 6-mer (scratch)       & 48.9\% & 0.989 & 0.171 \\
NT-v2 genus-balanced     & 37.2\% & 0.862 & 0.420 \\
\botrule
\end{tabular*}
\end{table}

\begin{figure*}[t]
\centering
\includegraphics[width=\textwidth]{cross_setting_comparison.pdf}
\caption{\textbf{Sample-level abundance fidelity across settings} on the
natural-distribution pool (9$\times$10K reads), sorted by Pearson $r$; read-level
accuracy is shown in parentheses. Abundance correlation stays high
($r\approx0.99$) down to $\sim$60\% read accuracy because independent per-read
errors cancel in aggregate. Genus-balancing is the clear outlier
($r=0.86$, Bray--Curtis $0.42$): it skews predictions toward rare genera and
distorts the natural composition.}
\label{fig:abund}
\end{figure*}

\subsection{Comparison with Kraken\,2: abundance versus detection}

On the coverage-matched genus pool (85{,}819 reads whose source species are in
the Kraken\,2 database), we separate two questions---\emph{how much} (abundance)
and \emph{whether present} (detection)---and compare the neural models against
Kraken\,2 both as raw report counts and with Bracken re-estimation
(Table~\ref{tab:kraken}, Figure~\ref{fig:kraken}). For abundance, raw Kraken\,2
counts give only Pearson $r=0.823$: Kraken\,2 leaves $\sim$30\% of reads
unclassified and, because that abstention is spread across taxa, it uniformly
deflates predicted abundances. This is a fixable artefact, not a fundamental
limit---Bracken \citep{Lu2017}, which re-estimates abundance from the same
Kraken\,2 report, raises genus abundance to $r=0.997$ (Bray--Curtis $0.060$) on
the same reads and database. Bracken therefore reaches the same near-perfect
range as the neural models (MT 13-mer $r=0.999$, NT-v2 $r=0.992$):
\textbf{the neural models have no sample-level abundance advantage over a
properly configured classical pipeline.} (Bracken redistributes the whole report and cannot be restricted per
read to the in-database subset; on the full 100K sample raw Kraken\,2 is
$r=0.845$, so this is a like-for-like comparison.) For detection, Kraken\,2---and
Bracken, which inherits its classifications---dominates (ROC AUC 0.966, 93.5\%
sensitivity at 95\% specificity), far above the 6-mer neural detectors (NT-v2
AUC 0.681, 17.5\%), with the overlapping-13-mer model in between (AUC 0.905,
66.5\%). On in-database data, then, the classical Kraken\,2$+$Bracken pipeline is
the stronger tool for \emph{both} abundance and detection; the neural models'
genuine advantages lie elsewhere---higher read-level accuracy under a good
tokenizer (MT 13-mer 94.9\% vs Kraken\,2 77.7\%) and graceful degradation
outside the reference database (below, where Kraken\,2 scores 0\% on 219
out-of-database species).

\begin{table*}[t]
\centering
\caption{\textbf{Genus-level, coverage-matched comparison with Kraken\,2}
(85{,}819 in-database reads). Raw Kraken\,2 abundance is deflated by its
$\sim$30\% abstention ($r=0.823$); re-estimating with Bracken raises Kraken\,2
genus abundance to $r=0.997$ (main text), on par with the best neural model, so
the neural advantage on abundance disappears against the proper pipeline.
Kraken\,2 also dominates detection. Sens@95 = sensitivity at 95\% specificity.}
\label{tab:kraken}
\begin{tabular*}{\textwidth}{@{\extracolsep{\fill}}lccccc@{}}
\toprule
Model & Read acc.\ & Pearson $r$ & Bray--Curtis & ROC AUC & Sens@95 \\
\midrule
MT 13-mer            & 94.9\% & \textbf{0.999} & 0.029 & 0.905 & 66.5\% \\
NT-v2 6-mer          & 68.9\% & 0.992 & 0.116 & 0.681 & 17.5\% \\
Kraken\,2 ($k{=}35$) & 77.7\% & 0.823 & 0.126 & \textbf{0.966} & \textbf{93.5\%} \\
MT 6-mer             & 51.5\% & 0.983 & 0.177 & 0.576 & 9.2\% \\
\botrule
\end{tabular*}
\end{table*}

\begin{figure}[t]
\centering
\includegraphics[width=\columnwidth]{tradeoff_abundance_detection.pdf}
\caption{\textbf{Abundance versus detection on the in-database pool}
(coverage-matched, 85{,}819 reads); marker area scales with ROC AUC. The neural
models reach high abundance fidelity (Pearson $r$) but weak detection
sensitivity. Raw Kraken\,2 is the opposite (strong detector, deflated abundance),
but \emph{Bracken re-estimation moves the classical pipeline to the top-right}
($r=0.997$, sensitivity 93.5\%)---best at both. On in-database data the neural
models therefore hold no advantage here; their value lies off this plot, in
read-level accuracy and out-of-database robustness (Section~\ref{sec:sample}).}
\label{fig:kraken}
\end{figure}

\subsection{Species level reveals the same tokenizer bottleneck}

At the species level (1{,}535 classes) the gap widens in Kraken\,2's favour on
the coverage-matched pool (77.2\% read accuracy, $r=0.847$), because exact long
$k$-mers resolve fine taxonomy that 6-mer neural models cannot: a flat NT-v2
6-mer species classifier reaches only 17.6\% Top-1, versus $\sim$50\% for the
overlapping-13-mer model. The natural remedy is a hierarchical pipeline---route a
read to its genus, then to a species within that genus---but this only helps with
a near-perfect router. Under \emph{oracle} genus routing, per-genus classifiers
lift species Top-1 to 29.5\%; with a \emph{realistic} learned router (genus
accuracy $\sim$66\%, itself the 6-mer ceiling of Section~\ref{sec:decomp}) the
same pipeline falls to 15.6\%---at or below the flat model. The bottleneck is
therefore the genus router, not the per-genus heads (oracle routing makes
genus-level accuracy $1.00$ by construction), so the 6-mer genus ceiling
propagates directly into species performance. Overlapping long-$k$ tokenization
is thus the shared fix for both genus routing and intra-genus separation; adding
a sub-genus split \emph{without} a better representation instead hurts (a
negative result; Supplementary). Conversely, Kraken\,2's dependence
on database coverage is a real liability---on 219 species outside its custom
database it scores 0\% while the neural models retain non-zero posterior mass,
the flip side of the abstention behaviour above.

Finally, the abstention-free behaviour of neural models has a cost for rare
taxa: a genuinely absent genus still receives a small number of spurious reads
(on the order of $\sim$0.28\% of a sample), setting a noise floor below which
detection is unreliable. This noise floor has a clean interpretation and
motivates the calibrated-threshold recommendation in
Section~\ref{sec:recommendations}.

\subsection{A real mock community: the in-database dominance does not transfer}
\label{sec:mock}

All comparisons so far use simulated reads. As a first check on real data, we
ran the trained genus classifiers and the Kraken\,2$+$Bracken pipeline
(inference only) on two Illumina WGS replicates of the ZymoBIOMICS Gut
Microbiome Standard (D6331; SRR33710519 and SRR33710518; 3M R1 reads each,
median 135\,bp). The vendor genomic-DNA percentages are the sample-level ground
truth. Of the 21 mock taxa, 11 genera (81.5\% of the community) fall within our
120 training genera; the remainder are out-of-set (dominated by
\emph{Escherichia} at 14\%, absent from the 120 genera), an unavoidable ceiling
for every method. Metrics are means over the two replicates
(Table~\ref{tab:mock}); per-replicate values and per-genus residuals are in the
Supplementary Material.

Two things change relative to the simulated in-database pool. First,
\emph{every} method's genus-abundance correlation collapses from
$r\approx0.99$--$0.997$ to $r\approx0.41$--$0.58$ (Spearman
$\rho\approx0.34$--$0.57$): the closed-set numbers do not transfer to real
reads. With only 11 in-set genera as points, bootstrap 95\% intervals on Pearson
$r$ are wide and overlapping, so we do \emph{not} rank methods by $r$.
Restricting to the nine genera with expected abundance $\ge$1\% raises all three
correlations into a tighter band ($r\approx0.57$--$0.62$), confirming that the
headline gap is driven partly by near-zero taxa and a few large residuals rather
than a stable method ordering. Second, the classical pipeline's simulated
dominance evaporates. Kraken\,2 classifies only $\sim$44\% of the real reads
($\sim$56\% unclassified, versus $\sim$30\% on simulation); even after Bracken,
detection sensitivity among genera expected at $\ge$1\% is 55.6\%---below the
neural models (NT-v2 83.3\%, MT 77.8\%)---and Bray--Curtis error is the worst of
the three (0.434), inflated by over-estimating \emph{Veillonella}. Always-labelling
neural models are competitive on composition error (MT best at 0.353; NT-v2
0.384).

Shared failure modes appear in both replicates and both neural models:
\emph{Roseburia} (17\% expected after in-set renormalization) is nearly missed
($\sim$0.5\%), while \emph{Clostridium} absorbs substantial mass
($\sim$12--18\% predicted vs $\sim$0 expected)---partly because
\emph{Clostridioides difficile} (1.5\% of the mock; out-of-set under current
taxonomy) may map into the training \emph{Clostridium} label if older names
were used. Out-of-set \emph{Escherichia} leaks into neighbouring
Enterobacteriaceae false positives. These patterns mix real-error behaviour with
strain/taxonomy mismatch, so the probe is not a pure sequencing-artefact test.

\begin{table}[t]
\centering
\caption{\textbf{Real mock community} (ZymoBIOMICS D6331; mean of two
replicates; 11 in-set genera, 18.5\% out-of-set). Pearson $r$ intervals overlap
heavily ($n{=}11$); Sens.\ = detection among genera with expected
abundance $\ge$1\% (predicted fraction $\ge$1\%). No method wins on all axes.}
\label{tab:mock}
\begin{tabular*}{\columnwidth}{@{\extracolsep{\fill}}lcccc@{}}
\toprule
Method & Pearson $r$ & Bray--Curtis & Sens.\ @1\% & Reads clf.\ \\
\midrule
Kraken\,2$+$Bracken & 0.578 & 0.434 & 55.6\% & 43.9\% \\
MT 13-mer           & 0.538 & 0.353 & 77.8\% & 100\% \\
NT-v2 6-mer         & 0.412 & 0.384 & 83.3\% & 100\% \\
\botrule
\end{tabular*}
\end{table}

This two-replicate probe (11 in-set genera, 135\,bp reads) remains preliminary
rather than a definitive multi-community benchmark. It does make one point
cleanly: the classical pipeline's crushing in-database advantage
(Section~\ref{sec:sample}) is a property of simulated, fully-covered data, not
of real communities---on real reads the neural models are competitive, and the
binding open problem for \emph{all} methods is the large gap between closed-set
and real-world accuracy. A fuller evaluation across additional communities,
matched read lengths, and further baselines is the priority follow-up
(Section~\ref{sec:limitations}).
